\subsubsection*{A one-derivative-loss pressure space}
\label{subsec:shifted-pressure-repair}
The isolated delayed-pressure primitive fails to preserve the
original doubled analytic coefficient norm, as demonstrated above.
The pressure-to-velocity composition nevertheless admits an exact
repair: retain the velocity weight $w_{a,b}$ of (37), but place
the pressure in the weaker norm
\begin{equation}\label{eq:shifted-P-norm}
 \|P\|_{\mathcal P}:=\sup_{a,b\ge0}\sup_{\eta\in[-1,1]}
 \frac{|\partial_\eta^bP_a(\eta)|}{w_{a,b+1}},
 \qquad P=\sum_{a\ge0}P_a(\eta)X^a .
\end{equation}
The additional derivative index is essential: the counterexample
has an order-four pole at radial degree two, while the shifted
weight allows its third-degree polynomial derivative growth.

\begin{lemma}[Sharp three-jet compensation]\label{lem:pressure-three-jet}
For all integers $a,b\ge0$ and $R,\rho>0$,
\begin{equation}\label{eq:shifted-exact-ratio}
 \frac{w_{a,b+3}}{(a+1)^2w_{a+1,b+1}}
 =\frac R{\rho^2}
 \frac{(2a+1)(2a+2)(a+2)^2}{(a+1)^4}
 \frac{(b+2)^2}{(b+4)^2}
 \le\frac{8R}{\rho^2}.
\end{equation}
The constant eight is sharp. More generally, if $n\ge2$,
$a\ge n-2$, and $j=0,1,2,3$, then
\begin{equation}\label{eq:shifted-four-slot}
 \frac{(a+n+1)^{3-j}w_{a,b+j}}
 {(a+1)^2w_{a+1,b+1}}
 \le32\,3^{3-j}R\rho^{1-j}.
\end{equation}
\end{lemma}
\begin{proof}
For $0\le j\le3$, the ratio without $(a+n+1)^{3-j}$
equals
\[
 R\rho^{1-j}\frac{(2a+1)(2a+2)(a+2)^2}{(a+1)^4}
 \frac{(b+2)^2}{(b+j+1)^2}
 \prod_{\ell=j+1}^{3}(2a+b+\ell)^{-1}.
\]
The radial numerator is at most eight times its denominator,
because
\[
8(a+1)^4-(2a+1)(2a+2)(a+2)^2
=2a^2(a+1)(2a+3)\ge0.
\]
The case $a=0,b\to\infty$ gives sharpness for $j=3$.
For all four jets use $(b+2)/(b+j+1)\le2$,
$2a+b+\ell\ge a+1$ and $a+n+1\le3(a+1)$.
\end{proof}
The norm also admits multiplication by an analytic function
$g(\eta)$ with $|\partial_\eta^\ell g|\le M_g\ell!$:
\begin{equation}\label{eq:P-multiplier}
 \|gP\|_{\mathcal P}\le\frac{80}{27}M_g\|P\|_{\mathcal P},
 \qquad 0<\rho\le1/4.
\end{equation}
Indeed the $\ell$-th Leibniz term has normalized factor at most
$\rho^\ell(\ell+1)^2$, whose sum is
$(1+\rho)/(1-\rho)^3\le80/27$.

\begin{theorem}[Complete isolated delayed-pressure primitive]
\label{thm:shifted-delayed-pressure}
Let $0<h<1/100$, $R>0$, $0<\rho\le1/4$, $n\ge3$ and
$F=U_{n-2}\in I_{n-2}$. Let $\Pi_n^{\rm del}$ be the pressure
primitive produced solely by the
$-Z_{0,n-2}^{[2]}V_{n-2}$ term of (28), with $V_{n-2}$
given by (24) and zero axis integration constant. Then
$\Pi_n^{\rm del}\in\mathcal P\cap I_{n-1}$ and
\begin{equation}\label{eq:shifted-delayed-bound}
 \|\Pi_n^{\rm del}\|_{\mathcal P}
 \le C_{\rm P,del}\|F\|_{\mathcal B},\qquad
 C_{\rm P,del}=\frac{320000000\,R}{27}
 \sum_{j=0}^{3}3^{3-j}\rho^{1-j}.
\end{equation}
\end{theorem}
\begin{proof}
Put $r=a+1$, $q=\lambda_{n-2}$,
$B=2\eta/L$, $G=(1-\eta^2)/L$, and
$Z_s^{(r)}=G\partial_\eta+B(s-r)$. The exact coefficient is
\[
 (\Pi_n^{\rm del})_{a+1}
 =\frac1{2r^2}Z_{q-D}^{(r)}Z_q^{(r)}
 [B(r-D-q)F_a-G\partial_\eta F_a].
\]
Its normal form has four channels,
$\sum_{j=0}^3c_{a,n,j}\partial_\eta^jF_a$.
Writing $t=a+n+1$, on $|\eta|\le3$ one has
$|B|<8$, $|G|<13$ and shifts at most $2t$.
The two $Z$ compositions, with two successive analytic
radius losses of $1/4$, increase the initial coefficient bound
by at most $100$ each. Thus on $|\eta|\le5/2$,
$|c_{a,n,j}|\le250000t^{3-j}$.
On $[-1,1]$ Cauchy's estimate gives
$|\partial_\eta^\ell c_{a,n,j}|
\le250000t^{3-j}\ell!$.
Combine the shifted four-jet estimate
\eqref{eq:shifted-four-slot} with the unsigned Leibniz
factor in \eqref{eq:P-multiplier}; the extra factor
$1/(2r^2)$ produces the displayed constant.
\end{proof}

\begin{theorem}[Bounded pressure-to-velocity reconstruction]
\label{thm:pressure-to-velocity}
Let $0<h<1/100$, $0<\rho\le1/4$, $R>0$, $n\ge1$.
For $P\in\mathcal P\cap I_{n-1}$, the actual axial
momentum pressure term belongs to the original velocity space
with valuation $I_n$ and
\begin{equation}\label{eq:P-to-U}
 \left\|\frac12J_1Z_{-2A,n}P\right\|_{\mathcal B}
 \le\frac{40}{27}
 \left(\frac{104R}{\rho^2}+\frac{168R}{\rho}\right)
 \|P\|_{\mathcal P}.
\end{equation}
\end{theorem}
\begin{proof}
At input degree $a\ge n-1$, put
$q=\lambda_n-2A=2h(n-1)-1$.
Then $-(a+1)\le q-a\le0$. Use the exact commutator
\[
 Z_{-2A,n}P
 =G\partial_\eta P+B(q-D_X)P
 =\partial_\eta(GP)-G'P+B(q-D_X)P.
\]
The original inverse weight-ratio bounds give
$\|J_1\partial_\eta H\|_{\mathcal B}
\le8R\rho^{-2}\|H\|_{\mathcal P}$ and
$\|J_1 H\|_{\mathcal B}
\le8R\rho^{-1}\|H\|_{\mathcal P}$.
The factor $|q-a|\le a+1$ cancels the residual inverse
denominator in the last term. Cauchy bounds on $|\eta|\le3$
give multiplier constants $M_G\le13$, $M_{G'}\le13$
and $M_B\le8$. Apply \eqref{eq:P-multiplier} and
sum the three contributions with prefactor $1/2$.
\end{proof}

The pressure counterexample is consistent with both theorems.
For $h=1/200$, $\rho=1/4$, $F=X/(1+4i\eta)$,
$\Pi_3^{\rm del}$ has a fourth-order pole at $\eta=i/4$
and belongs to $\mathcal P$ but not $\mathcal B$.
After applying $\frac12J_1Z_{-2A,3}$, its radial-degree-three
coefficient has a fifth-order pole with exact nonzero
coefficient $1044012500000i/19709998099203$. Its derivative
growth is admissible in the original radial-degree-three
velocity weight.

\paragraph{The corrected majorant and physical reconstruction.}
Set $B_n=\|\varphi_n\|_{I_n}+\|U_n\|_{I_n}$.
The exact coefficient identity $(I-K_n)w_n=S_nf_n$ has the
following rigorous \emph{conditional} consequence. If
$\|K_n\|\le C/(n+1)$ and the \emph{complete} regularized
forcing satisfies
\[
\|S_nf_n\|\le c_0(B_{n-1}+B_{n-2})
+c_1\sum_{i=1}^{n-1}B_iB_{n-i},
\]
then, for $C/(n+1)\le1/2$,
\begin{equation}\label{eq:corrected-velocity-recurrence}
B_n\le2c_0(B_{n-1}+B_{n-2})
+2c_1\sum_{i=1}^{n-1}B_iB_{n-i}.
\end{equation}
This velocity-only recurrence no longer assigns the pressure
primitive the stronger velocity norm, but the displayed
\emph{full-source} premise has yet to be verified against
every term of (25)--(28).

For pressure reconstruction, if
$\Pi_n\in\mathcal P\cap I_{n-1}$ and
$\|\Pi_n\|_{\mathcal P}\le MT^n$, then for $|X|\le R/8$
and each fixed integer $b\ge0$,
\begin{equation}\label{eq:P-derivative-sum}
|\partial_\eta^b\Pi_n(X,\eta)|
\le 2^{2-n}MT^n\frac{(b+1)!}{(b+2)^2}
\left(\frac2\rho\right)^{b+1}.
\end{equation}
The binomial bound
$\binom{2a+b+1}{b+1}\le4^a2^{b+1}$
and a geometric radial tail prove this directly.
Thus the pressure series and all fixed parameter derivatives
converge normally if $|s|T<2$.
A global geometric $\mathcal P$-bound for the \emph{complete}
pressure sequence, the actual full forcing estimate and the
termwise physical-coordinate PDE reconstruction remain
separate mathematical tasks. Neither Theorem~1.1 nor a
whole-space unforced singularity follows from this repair alone.
